\documentclass[11pt]{article}

\usepackage[a4paper,margin=2.5cm]{geometry}
\usepackage{amsmath,amssymb,amsthm}
\usepackage[colorlinks=true,linkcolor=blue,citecolor=blue,urlcolor=blue]{hyperref}
\usepackage{booktabs}
\usepackage{enumitem}
\usepackage{natbib}

\newtheorem{theorem}{Theorem}[section]
\newtheorem{lemma}{Lemma}[section]
\newtheorem{proposition}{Proposition}[section]
\newtheorem{corollary}{Corollary}[section]
\newtheorem{definition}{Definition}[section]
\newtheorem{remark}{Remark}[section]

\newcommand{\twoI}{2I}
\newcommand{\Vsixhundred}{V_{600}}
\newcommand{\Dicfive}{\mathrm{Dic}_5}
\newcommand{\Vtwentyfour}{V_{24}}
\newcommand{\sigmaG}{\sigma}
\newcommand{\Eq}{E_q}
\newcommand{\Tcasc}{T_{\mathrm{casc}}}

\title{A Discrete Hawking Quantum from $\sigma$-Pair Excitations\\
       on the Binary Icosahedral Group}
\author{%
  Lee Smart \\[3pt]
  \textit{Institute of Vibrational Field Dynamics} \\[2pt]
  \texttt{contact@vibrationalfielddynamics.org}}
\date{May 2026}

\begin{document}

\maketitle

\noindent\textbf{Status:} Preprint (not peer-reviewed).

\begin{abstract}
We define the $\sigma$-Galois pair-excitation observable
$E(v) := |v - \sigmaG(v)|^2_{\mathrm{Eucl}}$ on the binary icosahedral
group $\twoI = \Vsixhundred$, where $\sigmaG : \sqrt{5} \mapsto -\sqrt{5}$
acts componentwise on icosian quaternions. We prove that
$E(v) = 0$ for the $24$ $\sigmaG$-fixed vertices (the $24$-cell sublattice
$\Vtwentyfour$) and $E(v) = 5/2$ \emph{exactly} for each of the
remaining $96$ $\sigmaG$-mobile vertices. The numerical value $5/2$
follows from a direct coordinate calculation; constancy of $E$ on the
$\sigmaG$-mobile sector additionally follows from a symmetry-grade
argument: the bilateral action of $\Vtwentyfour$,
$v \mapsto g v h$ for $g, h \in \Vtwentyfour$, is transitive on
$\sigmaG$-mobile vertices and commutes with $\sigmaG$.
Combined with the Paper~1~\citep{V600Paper1} per-coset incidence
$|gH \cap \Vtwentyfour|/|gH \setminus \Vtwentyfour| = 4/16 = 1/4$ for
$H = \Dicfive$, this yields an exact discrete first-law identity per
coset (left or right, bulk or non-bulk):
\(
  E_C = \Tcasc \cdot A_C, \quad S_C = A_C/4, \quad
  \Tcasc \cdot S_C = E_C / 4,
\)
with $\Tcasc = \Eq = 5/2$ in Euclidean cascade units and $A_C = 16$,
$S_C = 4$, $E_C = 40$. Interpreting $\sigmaG$-mobile vertices as Hawking radiation
channels and $\sigmaG$-fixed vertices as horizon anchor states gives a
discrete monochromatic analogue of Hawking radiation. The semiclassical
recovery of a Planck continuum from the discrete gap and the calibration
of $\Eq$ to $\hbar c^3 / (8 \pi G M k_B)$ are stated as Tier-2 work.
\end{abstract}

\section{Introduction}\label{sec:intro}

The Hawking temperature $T_H = \hbar c^3/(8 \pi G M k_B)$
\citep{Hawking1975}, in conjunction with the Bekenstein--Hawking entropy
$S_{\mathrm{BH}} = A/4$ in natural units \citep{Bekenstein1973}, defines
a thermodynamic structure for black-hole horizons. In the semiclassical
formulation the radiation spectrum is approximately Planckian: the
horizon emits a continuum of modes with thermal occupation at temperature
$T_H$. The continuum nature is fundamental in QFT in curved spacetime.

This paper reports a structurally distinct picture: in the finite
arithmetic of the binary icosahedral group $\twoI$ — equivalently the
$120$ vertices of the $600$-cell $\Vsixhundred$ — the natural
$\sigmaG$-pair excitation observable yields a single quantum gap
$\Eq = 5/2$ shared by all $\sigmaG$-mobile vertices, not a continuum.
Constancy of the spectrum on the $\sigmaG$-mobile sector is explained
by a transitive symmetry — the bilateral action of the
$\sigmaG$-fixed sublattice $\Vtwentyfour$ on $\Vsixhundred$ — and the
numerical value is computed in coordinates.
Combined with the Paper~1~\citep{V600Paper1} coset-incidence identity
$S/A = 1/4$, this produces an exact discrete first-law analogue
$\Tcasc \cdot S = E/4$ with $\Tcasc = \Eq$.

The paper does not derive physical Hawking radiation. It claims a
finite-arithmetic identity in $\Vsixhundred$ that has the algebraic
shape of a discrete Hawking quantum, with the semiclassical recovery
of the Planck continuum and the calibration to physical units stated
as future work. Section~\ref{sec:setup} fixes notation and imports the
$\Vsixhundred$/$\Dicfive$/$\Vtwentyfour$ structure from Paper~1.
Section~\ref{sec:obs} defines the $\sigmaG$-pair excitation observable.
Section~\ref{sec:main} states and proves the monochromatic-spectrum
theorem. Section~\ref{sec:firstlaw} derives the per-coset first-law
identity. Section~\ref{sec:scope} sets the scope.

\section{Setup}\label{sec:setup}

We import the following from Paper~1~\citep{V600Paper1}; the present
paper does not re-prove any of these statements.

\begin{itemize}[leftmargin=2em]
\item $K := \mathbb{Q}(\sqrt{5})$ with Galois twist
$\sigmaG : a + b\sqrt{5} \mapsto a - b\sqrt{5}$, extended componentwise
to the icosian quaternion algebra $\mathbb{H}_K$.
\item $\Vsixhundred = \twoI \subset \mathrm{Sp}(1, K)$: the $120$ unit
icosians forming the binary icosahedral group, partitioned into
$8$~axis, $16$~half-type, and $96$~even-permutation vertices.
\item $\Vtwentyfour := \{ v \in \Vsixhundred : \sigmaG(v) = v \}$: the
$\sigmaG$-fixed sublattice, equal to the $8$ axis vertices plus the $16$
half-type vertices and forming the regular $24$-cell. This is also
$\Vsixhundred \cap \sigmaG(\Vsixhundred)$.
\item $\sigmaG$-\emph{mobile} vertices: $\Vsixhundred \setminus \Vtwentyfour$,
the $96$ even-permutation vertices.
\item $H := \Dicfive \subset \twoI$: the bulk subgroup of order $20$,
the union of the $K=72$ and $K=0$ cycles of $T_\tau$. $H$ is
self-normalizing in $\twoI$.
\item Per-coset incidence (Paper~1, Theorem on coset incidence):
for every left coset $gH$ and right coset $Hg$,
$|gH \cap \Vtwentyfour| = |Hg \cap \Vtwentyfour| = 4$ and
$|gH \setminus \Vtwentyfour| = |Hg \setminus \Vtwentyfour| = 16$.
\end{itemize}

\section{The $\sigmaG$-pair excitation}\label{sec:obs}

\begin{definition}[$\sigmaG$-pair excitation]\label{def:Eobs}
For each vertex $v \in \Vsixhundred$, define
\[
  E(v) \;:=\; |v - \sigmaG(v)|^2_{\mathrm{Eucl}},
\]
the squared icosian-Euclidean distance between $v$ and its $\sigmaG$-image,
computed exactly in $K = \mathbb{Q}(\sqrt{5})$. The value lies in
$\mathbb{Q}$ because the $\sqrt{5}$-component of $|v - \sigmaG(v)|^2$
vanishes by $\sigmaG$-equivariance of the squared norm.
\end{definition}

\begin{lemma}[Vanishing on $\Vtwentyfour$]\label{lem:vanish}
$E(v) = 0$ if and only if $v \in \Vtwentyfour$.
\end{lemma}

\begin{proof}
$E(v) = 0$ iff $v - \sigmaG(v) = 0$ iff $\sigmaG(v) = v$ iff
$v \in \Vtwentyfour$ by Definition of the $24$-cell.
\end{proof}

\begin{lemma}[Coordinate identity on $\sigmaG$-mobile vertices]\label{lem:coord}
For every $\sigmaG$-mobile vertex $v \in \Vsixhundred \setminus \Vtwentyfour$,
the difference $v - \sigmaG(v)$ has exactly \emph{two} nonzero coordinates,
each of value $\pm \sqrt{5}/2 \in K$.
\end{lemma}

\begin{proof}
The $\sigmaG$-mobile vertices are precisely the $96$ even-permutation
vertices, with coordinate orbit $(0,\, \pm 1/(2\varphi),\, \pm 1/2,\,
\pm \varphi/2)$ up to even permutation. Recall the $K$-arithmetic
identities $1/(2\varphi) = (\sqrt{5}-1)/4$ and
$\varphi/2 = (1+\sqrt{5})/4$. Applied componentwise:
\begin{align*}
\sigmaG(0) &= 0, \\
\sigmaG(1/2) &= 1/2, \\
\sigmaG\bigl(1/(2\varphi)\bigr) &= -(1+\sqrt{5})/4 = -\varphi/2, \\
\sigmaG\bigl(\varphi/2\bigr) &= (1-\sqrt{5})/4 = -1/(2\varphi),
\end{align*}
so $\sigmaG$ sends the golden coordinate values $1/(2\varphi)$ and
$\varphi/2$ to their conjugate partners ($-\varphi/2$ and $-1/(2\varphi)$
respectively) in the same coordinate slots, while leaving $0$ and
$\pm 1/2$ slots invariant.
Subtracting,
\[
  1/(2\varphi) - \sigmaG(1/(2\varphi)) =
  (\sqrt{5}-1)/4 + (1+\sqrt{5})/4 = \sqrt{5}/2,
\]
\[
  \varphi/2 - \sigmaG(\varphi/2) =
  (1+\sqrt{5})/4 - (1-\sqrt{5})/4 = \sqrt{5}/2,
\]
and the $0$ and $\pm 1/2$ coordinates contribute zero to $v - \sigmaG(v)$.
So exactly two coordinates of $v - \sigmaG(v)$ are nonzero, each of
value $\pm \sqrt{5}/2$ depending on the original sign convention.
\end{proof}

\section{Main theorem: monochromatic spectrum}\label{sec:main}

\begin{theorem}[Hawking quantum]\label{thm:Eq}
The $\sigmaG$-pair excitation observable $E$ takes exactly two values on
$\Vsixhundred$: $E(v) = 0$ for all $24$ vertices in $\Vtwentyfour$ and
\[
  E(v) \;=\; \Eq \;:=\; \tfrac{5}{2}
\]
for each of the remaining $96$ $\sigmaG$-mobile vertices. The spectrum
is the discrete distribution
\(
  24\,\delta_{0} + 96\,\delta_{5/2}.
\)
\end{theorem}

\begin{proof}
Vanishing on $\Vtwentyfour$ is Lemma~\ref{lem:vanish}.

For $\sigmaG$-mobile $v$, Lemma~\ref{lem:coord} gives $v - \sigmaG(v)$
having exactly two nonzero coordinates each of value $\pm \sqrt{5}/2$.
Computing the squared Euclidean norm:
\[
  E(v) \;=\; |v - \sigmaG(v)|^2 \;=\; (\sqrt{5}/2)^2 + (\sqrt{5}/2)^2
  \;=\; \tfrac{5}{4} + \tfrac{5}{4} \;=\; \tfrac{5}{2}.
\]
This holds for every $\sigmaG$-mobile vertex regardless of which two
coordinate positions are nonzero, so $E(v) = 5/2$ uniformly across all
$96$ such vertices.

Symmetry corroboration. Define the bilateral action $\Gamma$ on
$\Vsixhundred$ by $\Gamma := \{ v \mapsto g v h : g, h \in \Vtwentyfour \}$.
Each generator is a quaternionic isometry (left- and right-multiplication
by unit quaternions preserve $|\cdot|^2$) and commutes with $\sigmaG$
because $g, h \in \Vtwentyfour$ have $\sigmaG(g) = g$ and $\sigmaG(h) = h$,
so $\sigmaG(g v h) = g \sigmaG(v) h$. Therefore $E$ is constant on
$\Gamma$-orbits. The $\Gamma$-orbits on $\Vsixhundred$ are the
$\Vtwentyfour$-double cosets $\Vtwentyfour \cdot v \cdot \Vtwentyfour$;
the certificate \texttt{verify.py} confirms that there are exactly two
such orbits — $\Vtwentyfour$ itself (of size $24$) and a single orbit of
size $96$ (the $\sigmaG$-mobile sector). Constancy of $E$ on the
$96$-orbit is therefore implied by $\Gamma$-invariance.
\end{proof}

We refer to $\Eq = 5/2$ as the \emph{cascade Hawking quantum}.

\section{First-law identity}\label{sec:firstlaw}

\begin{definition}[Coset thermodynamic content]\label{def:coset_thermo}
For a left coset $C = gH$ of $H = \Dicfive$ in $\twoI$, set
\[
  S(C) := |C \cap \Vtwentyfour|, \quad
  A(C) := |C \setminus \Vtwentyfour|, \quad
  E(C) := \!\sum_{v \in C \setminus \Vtwentyfour}\! E(v),
  \quad \Tcasc(C) := E(C) / A(C).
\]
The same definitions apply to right cosets $Hg$.
\end{definition}

\begin{theorem}[Discrete first law]\label{thm:firstlaw}
For every coset $C$ of $H$ in $\twoI$ — left or right, bulk or
non-bulk —
\[
  S(C) = 4, \quad A(C) = 16, \quad E(C) = 40, \quad
  \Tcasc(C) = \Eq = \tfrac{5}{2},
\]
and the following identities hold exactly:
\[
  E(C) = \Tcasc(C) \cdot A(C), \qquad
  S(C) = A(C) / 4, \qquad
  \Tcasc(C) \cdot S(C) = E(C) / 4.
\]
Aggregating over the five non-bulk left cosets:
\(
  S_{\mathrm{tot}} = 20,\; A_{\mathrm{tot}} = 80,\;
  E_{\mathrm{tot}} = 200,\;
  \Tcasc \cdot S_{\mathrm{tot}} = E_{\mathrm{tot}} / 4 = 50.
\)
\end{theorem}

\begin{proof}
$S(C) = 4$ and $A(C) = 16$ are the per-coset incidence counts of
Paper~1~\citep{V600Paper1}. By Theorem~\ref{thm:Eq}, every $\sigmaG$-mobile
vertex contributes $\Eq = 5/2$ to $E(C)$, so
\(
  E(C) = A(C) \cdot \Eq = 16 \cdot 5/2 = 40,
\)
and $\Tcasc(C) = E(C)/A(C) = 5/2$. The three identities are direct
substitutions:
$E(C) = \Tcasc(C) \cdot A(C)$ by definition of $\Tcasc(C)$;
$S(C) = 4 = 16/4 = A(C)/4$ by Paper~1;
$\Tcasc(C) \cdot S(C) = (5/2) \cdot 4 = 10 = 40/4 = E(C)/4$. The bulk
coset $C = H$ itself satisfies the same per-coset identity because
$|H \cap \Vtwentyfour| = 4$ (Paper 1, $24$-cell incidence with $H$).

The aggregate values follow by summing over the five disjoint non-bulk
cosets.
\end{proof}

\begin{remark}[Bekenstein--Hawking analogy, interpretive]\label{rem:BH}
Under the interpretive identification of $\sigmaG$-mobile vertices as
horizon channels (``area'') and $\sigmaG$-fixed vertices as anchor
states (``entropy'' bits), Theorem~\ref{thm:firstlaw} produces three
formal relations within this finite model that echo the structural
shape of the Bekenstein--Hawking entropy formula: a sum
$E = \Tcasc \cdot A$ that records $E$ as the channel count times a
common per-channel quantum, the ratio $S = A/4$ that mirrors the
Bekenstein coefficient (this part is genuinely standard, by
Paper~1~\citep{V600Paper1}), and a derived identity
$\Tcasc \cdot S = E/4$ that follows from those two by direct
substitution. None of these is the standard physical first law
$dE = T\,dS$ of black-hole thermodynamics; the present finite model
does not address dynamical changes. Calibration of $\Eq = 5/2$ to the
physical Hawking temperature
$T_H^{\mathrm{phys}} = \hbar c^3 / (8\pi G M k_B)$ is also not
constructed here.
\end{remark}

\section{Scope}\label{sec:scope}

\noindent
\fbox{\parbox{\dimexpr\linewidth-2\fboxsep-2\fboxrule\relax}{%
\textbf{Scope.}
\begin{itemize}
\item \textbf{Established (exact, unconditional):} the $\sigmaG$-pair
spectrum on $\Vsixhundred$ is $24\,\delta_0 + 96\,\delta_{5/2}$
(Theorem~\ref{thm:Eq}); the per-coset first-law identity
$\Tcasc \cdot S = E/4$ with $\Tcasc = \Eq = 5/2$, $A = 16$, $S = 4$,
$E = 40$ (Theorem~\ref{thm:firstlaw}); the bilateral $\Vtwentyfour$
symmetry orbit certificate confirms the $96$-orbit structure.
\item \textbf{Interpretive:} the identification of $\sigmaG$-mobile
vertices as Hawking radiation channels and $\sigmaG$-fixed vertices
as horizon anchor states (Remark~\ref{rem:BH}). A different
identification would attach different physical content to the same
arithmetic identity.
\item \textbf{Out of scope:} cascade-unit calibration to physical
Planck units; semiclassical recovery of a Planck continuum from the
discrete gap; quantitative primordial-black-hole evaporation cutoff;
the canonical involution $\tau_\sigma$ (Paper~3); cosmological
applications (Paper~4); ontology of $\Vsixhundred$. The formal
analogy between $\Eq$ and a Hawking-temperature parameter is
observed without claim; the unit calibration required to lift
$\Tcasc$ from a structural finite identity to a quantitative
physical Hawking temperature is not constructed here.
\end{itemize}}}

\section{Toward a calibrated continuum limit}\label{sec:pbh}

The finite spectrum derived here is qualitatively unlike a smooth
Planck continuum. To connect the present arithmetic identity to a
physical observable would require (i) a calibration that maps the
cascade unit $\Eq$ to a physical energy scale, (ii) a mechanism by
which sums over many cascade resolutions broaden the discrete gap
into an approximately thermal distribution, and (iii) an observational
model that distinguishes a discrete cutoff from a thermal tail in any
candidate setting (for instance, late-stage primordial black-hole
evaporation \citep{CarrEtAlPBH2021}). None of (i)--(iii) is
constructed here, and we make no quantitative observational claim. We
note the qualitative shape of the gap as a target for future
calibration work.

\section{Conclusion}\label{sec:concl}

The $\sigmaG$-Galois pair-excitation observable on the binary
icosahedral group $\twoI$ admits a closed-form spectrum:
$24$~zero modes (the $\sigmaG$-fixed $24$-cell sublattice) and
$96$~mobile modes at the single quantum value $\Eq = 5/2$
(Theorem~\ref{thm:Eq}). Combined with the Paper~1 coset-incidence
identity, this produces a per-coset first-law identity
$\Tcasc \cdot S = E/4$ with $\Tcasc = \Eq$ (Theorem~\ref{thm:firstlaw}). The
result is exact, parameter-free, and verifiable by the accompanying
script. We frame this as a discrete arithmetic analogue of the
Hawking-radiation thermodynamic structure, deferring physical-unit
calibration and semiclassical-continuum recovery to future work.

\appendix

\section{Verification certificate}\label{app:verify}

The companion script \texttt{verify.py} performs the following exact
checks and exits with status $0$ when all assertions hold:
\begin{enumerate}
\item $|\Vsixhundred| = 120$, $|\Vtwentyfour| = 24$,
$|\Vsixhundred \setminus \Vtwentyfour| = 96$.
\item For all $24$ $\sigmaG$-fixed vertices, $E(v) = 0$ exactly.
\item For all $96$ $\sigmaG$-mobile vertices, $E(v) = 5/2$ exactly,
with the rational value computed in $\mathbb{Q}(\sqrt{5})$.
\item Coordinate proof: for each of the $96$ $\sigmaG$-mobile
vertices, the difference $v - \sigmaG(v)$ has exactly two nonzero
coordinates with rational part $0$ and $\sqrt{5}$-part $\pm 1/2$.
\item Bilateral $\Vtwentyfour$ symmetry: $\sigmaG$ commutes with
left-multiplication by every $g \in \Vtwentyfour$ (24 generators), and
the bilateral action $\Gamma = \{v \mapsto g v h : g, h \in \Vtwentyfour\}$
on $\Vsixhundred$ has exactly two orbits — $\Vtwentyfour$ (size $24$)
and the $\sigmaG$-mobile sector (size $96$).
\item Per-coset first-law identity for all $6$ left cosets and all $6$
right cosets (bulk and non-bulk): $A=16$, $S=4$, $E=40$, $\Tcasc = 5/2$,
$\Tcasc \cdot S = E/4 = 10$.
\item Aggregate first law: $A_{\mathrm{tot}} = 80$, $S_{\mathrm{tot}} = 20$,
$E_{\mathrm{tot}} = 200$, $\Tcasc \cdot S_{\mathrm{tot}} = E_{\mathrm{tot}}/4 = 50$.
\end{enumerate}

The script depends only on the Python standard library and the shared
\texttt{vfd\_v600} package
(\texttt{papers/v600-programme/lib}). Reproducibility command:
\begin{verbatim}
PYTHONPATH=papers/v600-programme/lib \
    python3 papers/v600-hawking-quantum/verify.py
\end{verbatim}

\bibliographystyle{plainnat}
\bibliography{references}

\end{document}
